\section*{第 \ref{ch:b1:p-block-16-18} 章——p 区元素 II：氧、卤素与稀有气体}

\begin{solution}{exo:b1:p-block-16-18:1}
\ce{H2S} $-$II; \ce{S8} 0; \ce{SO2} $+$IV; \ce{SO3^2-} $+$IV; \ce{SO4^2-}
$+$VI；\ce{S2O3^2-} 中 S 的平均\omterm{def:b1:nernst:oxidation-number}{氧化数}为 $+$II；\ce{H2S2O7} $+$VI。
\end{solution}

\begin{solution}{exo:b1:p-block-16-18:2}
\omterm{def:b1:p-block-16-18:halogen}{卤素}能氧化位于其下方的卤化物：\ce{Cl2 + 2I- -> 2Cl- + I2} 及
\ce{Cl2 + 2Br- -> 2Cl- + Br2} 均可发生；碘与溴化物、溴与氯化物
之间则不能发生。
\end{solution}

\begin{solution}{exo:b1:p-block-16-18:3}
\ce{SO2}：S 有一对\omterm{def:b1:lewis-resonance-vsepr:lewis}{孤电子对}，在两个\omterm{def:b1:lewis-resonance-vsepr:resonance}{共振结构}中各含一个 S=O 与一个 \ce{S+-O-}
（或取两个 S=O 且带扩展八隅体）：角形。
\ce{SO3}：平面三角形，三个等价的 S--O 键。\ce{O3}：中心 O
原子含一对\omterm{def:b1:lewis-resonance-vsepr:lewis}{孤电子对}，两种\omterm{def:b1:lewis-resonance-vsepr:resonance}{共振结构}分别含一个 O=O 与一个 O--O$^-$：角形。
\end{solution}

\begin{solution}{exo:b1:p-block-16-18:4}
\ce{SO2 + H2O <=> H2SO3}; \ce{SO2 + 2NaOH -> Na2SO3 + H2O};
\ce{SO3 + H2O -> H2SO4}; \ce{SO3 + 2NaOH -> Na2SO4 + H2O}.
\end{solution}

\begin{solution}{exo:b1:p-block-16-18:5}
HF: $x^2/(0.10 - x) = 10^{-3.16}$, $x = \qty{8.0e-3}{mol/L}$, pH 2.10. HCl:
pH 1.00。H--F 键很强，氟离子强烈地形成\omterm{def:b1:intermolecular-forces-solvents:hbond}{氢键}，
故 HF 为\omterm{def:b1:predominant-reaction:strong-acid}{弱酸}。
% ledger: dfg:HF_ao, dfg:F-_ao
\end{solution}

\begin{solution}{exo:b1:p-block-16-18:6}
\ce{BrF5}：五个键和一对\omterm{def:b1:lewis-resonance-vsepr:lewis}{孤电子对}，四方锥形。\ce{IF7}：
五角双锥形。\ce{ICl4-}：四个键和两对\omterm{def:b1:lewis-resonance-vsepr:lewis}{孤电子对}，平面正
方形。\ce{XeO3}：三个键和一对\omterm{def:b1:lewis-resonance-vsepr:lewis}{孤电子对}，三角锥形。
\end{solution}

\begin{solution}{exo:b1:p-block-16-18:7}
$\log K = 2(1.36 - 0.53)/0.059 = 28$。析出的碘与淀粉形成深
蓝色\omterm{def:b1:complexation:complex}{配合物}。
\end{solution}

\begin{solution}{exo:b1:p-block-16-18:8}
$E^\circ(\ce{O3}/\ce{O2}) = 2.07 > 0.53$~V: \ce{O3 + 2I- + 2H+ -> O2 + I2
+ H2O}。生成的碘用硫代硫酸盐溶液\omterm{def:b1:titration-methods:titration}{滴定}
（见\cref{exo:b1:nernst:7}）：一个 \ce{I2} 对应一个 \ce{O3}。
\end{solution}

\begin{solution}{exo:b1:p-block-16-18:9}
\ce{C12H22O11 -> 12C + 11H2O}; $M = \qty{342.0}{g/mol}$, $10.0/342.0 =
\qty{0.0292}{mol}$，碳 $0.0292 \times 12 \times 12.0 = \qty{4.21}{g}$。
\end{solution}

\begin{solution}{exo:b1:p-block-16-18:10}
$x(0.10 + x)/(0.10 - x) = 10^{-1.99}$: $x = \qty{8.6e-3}{mol/L}$,
$[\ce{H3O+}] = 0.1086$，$\mathrm{pH} = 0.96$，而非 1.00：第二步
电离使 $[\ce{H3O+}]$ 约增加 \qty{9}{\percent}。
\end{solution}

\begin{solution}{exo:b1:p-block-16-18:11}
\ce{F2}/\ce{F-}（2.89~V）位于其他所有电对之上：没有任何\omterm{def:b1:nernst:couple}{氧化剂}能从
氟离子上夺取电子。水溶液电解时，水在低得多的电势下被
氧化（\ce{O2}/\ce{H2O}，1.23~V），会优先发生反应；
而熔融氟化物中不含水。
\end{solution}

\begin{solution}{exo:b1:p-block-16-18:12}
$\Delta_r G^\circ = -51.4 - (16.40 - 51.57) = \qty{-16.2}{kJ/mol}$,
$\log K = 16.2/5.708 = 2.84$, $K = 7 \times 10^2$.
% ledger: dfg:I2_ao, dfg:I-_ao, dfg:I3-_ao
\end{solution}

\begin{solution}{pb:b1:p-block-16-18:1}
\textbf{1.} 0, $+$IV, $+$VI, $+$VI.
\textbf{2.} \ce{S + O2 -> SO2}.
\textbf{3.} S 有两个\omterm{def:b1:lewis-resonance-vsepr:lewis}{成键电子对}和一对\omterm{def:b1:lewis-resonance-vsepr:lewis}{孤电子对}：角形，约
\ang{120}.
\textbf{4.} $10^6/32.1 = \qty{3.12e4}{mol}$.
\textbf{5.} \ce{O2} 为 \qty{3.115e4}{mol}：$V = 3.115 \times 10^4 \times
8.314 \times 298.15/10^5 = \qty{772}{m^3}$，所需空气 $772/0.21 = \qty{3.7e3}{m^3}$。
\textbf{6.} 燃烧放热很多；\omterm{def:b1:elementary-steps:catalyst}{催化剂}只在有限的温度范围内起作用，
而 \ce{SO2} 的氧化同样放热，
高温时转化不够完全。
\textbf{7.} \ce{2SO2 + O2 <=> 2SO3}.
\textbf{8.} 平衡常数很大并不说明\omterm{def:b1:rate-laws:rate}{反应速率}的大小：在
室温下反应并不进行。
\textbf{9.} 它提供了\omterm{def:b1:rate-laws:activation-energy}{活化能}较低的途径，且
本身可再生：属于\omterm{def:b1:elementary-steps:catalyst}{多相催化}剂。
\textbf{10.} 平面三角形。
\textbf{11.} 氧气越多，使商
$Q = p_{\ce{SO3}}^2/(p_{\ce{SO2}}^2 p_{\ce{O2}})$ 保持较低，
\ce{SO2} 的转化进行得更彻底。
\textbf{12.} \qty{0.3}{\percent}.
\textbf{13.} \ce{SO3 + H2SO4 -> H2S2O7}.
\textbf{14.} 它形成由细小酸滴组成的酸雾，穿过
吸收塔。
\textbf{15.} \ce{H2S2O7 + H2O -> 2H2SO4}.
\textbf{16.} \ce{2S + 3O2 + 2H2O -> 2H2SO4}.
\textbf{17.} 每个硫原子对应一个水分子：$\num{3.115e4} \times 18.0 =
\qty{561}{kg}$.
\textbf{18.} 稀释放出大量热量；若把水倒入酸中，
水会立即沸腾并使酸液飞溅。
\textbf{19.} $x(0.050 + x)/(0.050 - x) = 10^{-1.99}$: $x = \qty{7.5e-3}{mol/L}$,
$[\ce{H3O+}] = \qty{0.0575}{mol/L}$, $\mathrm{pH} = 1.24$.
\textbf{20.} $3.12 \times 10^4 \times 0.995 \times 98.1 = \qty{3.04}{t}$.
\textbf{21.} $84 \times 98.1/32.1 = \qty{257}{Mt}$.
\textbf{22.} $98.1/32.1 =$ \textbf{每吨硫可得 \qty{3.06}{t}
硫酸}。
\end{solution}
